% Generated draft: requires mathematical and source review.
\documentclass[11pt]{article}
\usepackage[margin=25mm]{geometry}
\usepackage{fontspec}
\setmainfont{Latin Modern Roman}
\usepackage{amsmath,amssymb,mathtools,unicode-math}
\setmathfont{Latin Modern Math}
\usepackage{enumitem,needspace,xurl,hyperref}
\setlength{\parindent}{0pt}
\setlength{\parskip}{5pt}
\setlength{\emergencystretch}{3em}
\newcommand{\N}{\mathbb{N}}
\newcommand{\Z}{\mathbb{Z}}
\newcommand{\Q}{\mathbb{Q}}
\newcommand{\R}{\mathbb{R}}
\newcommand{\C}{\mathbb{C}}
\newcommand{\F}{\mathbb{F}}
\newcommand{\A}{\mathbb{A}}
\newcommand{\PP}{\mathbb{P}}
\newcommand{\E}{\mathbb{E}}
\newcommand{\eps}{\varepsilon}
\newcommand{\dd}{\,\mathrm{d}}
\begin{document}
\section*{UnsolvedMath ID: 8800056}
\textbf{Source problem code:} catalogue entry
\subsection{Definitions and mathematical statement}
A pairing-based system typically uses efficiently represented finite groups and a computable nondegenerate bilinear map $e:G_1\times G_2\to G_T$, with
\[
e(aP,bQ)=e(P,Q)^{ab}
\]
for group elements and scalars. This supplies context, not a specification of the cryptographic scheme in the source.

Identify the particular system and parameters in the slides, and either find an attack achieving its intended break or establish resistance to the specified known attacks. A precise attack claim requires an adversarial task, access model, success criterion, and resource bound. A formal computational security formulation would measure adversarial advantage at a security parameter and distinguish nonnegligible efficient attacks from negligible advantage. The source does not give that game or the system's algorithms.

Resistance to a catalog of known attacks is different from an unconditional security proof against all attacks, and the source asks only the former alternative. Any reduction-based guarantee must identify its hardness assumption and parameter regime. The supplied summary mentions BLS, bilinear Diffie--Hellman-type assumptions, and no known generic break, but does not identify those systems with the particular slide system or supply a security theorem. No attack or blanket security claim for pairing cryptography is invented.
\subsection{Short English statement}
Analyze the specific pairing-based system from the missing slide context: find a defined break or assess its resistance to specified known attacks. The scheme, parameters, adversary model, and resource bounds are absent.
\subsection{Sources}
\begin{itemize}
\item Tanja Lange, IPAM \emph{Workshop on Number Theory and Cryptography --- Open Problems} (2006), attachment 6239, PDF page 41, printed slide 39, schedule source \url{https://www.ipam.ucla.edu/programs/workshops/workshop-i-number-theory-and-cryptography-open-problems/?tab=schedule}; accessed 2026-07-29. The supplied provenance records a PDF/legacy-PPT visual audit.
\item The research summary mentions BLS and BDH-based systems and related hardness assumptions without a particular scheme specification or publication citation.
\item Local registry \texttt{data/unsolvedmath/problems.json}, permanent UnsolvedMath ID $8800056$, source code catalogue entry, \emph{Breaking the pairing system}. Dataset attribution: Ulam AI and the UnsolvedMath Contributors.
\end{itemize}
\end{document}
